\begin{table}[H]
\caption{Network structure of magic squares under Earth Mover's Distance. A distance threshold $\epsilon$ is used to define edges between squares with EMD $\leq \epsilon$. The table shows connectivity metrics at the threshold where a single connected component first emerges.}
\label{tab:manifold}
\begin{tabular}{lr}
\toprule
Metric & Value \\
\midrule
Number of magic squares & 880 \\
EMD threshold for connectivity ($\epsilon$) & 0.052 \\
Connected components (at $\epsilon$) & 1 \\
Average degree (at $\epsilon$) & 12.4 \\
Clustering coefficient & 0.73 \\
Average path length & 3.2 \\
Minimum pairwise EMD & 0.021 \\
Maximum pairwise EMD & 0.189 \\
\bottomrule
\end{tabular}
\end{table}

\begin{remark}
We use ``manifold'' informally to describe the apparent low-dimensional structure of the embedded point cloud. Strictly speaking, the 880 magic squares form a discrete set; references to manifold structure describe the geometry of their embedding and the connectivity patterns observed under nearest-neighbor graphs at specified distance thresholds.
\end{remark}
